\section{System Model and Problem Formulation}\label{sec:model}

\paragraph{Dynamic graph and workload.}
The input is a stream of windows $t=1,\dots,T$. Window $t$ inserts and deletes edges and
vertices, which yields the graph $G_t=(V_t,E_t)$, and it marks a set of \emph{touched} vertices
(endpoints of new interactions). As in incremental GNN inference, the active set $A_t$ (the
touched vertices and their neighbours) is recomputed in every window by $S$ supersteps of a
GraphSAGE-style layer
\begin{equation}\label{eq:layer}
  h_v \leftarrow \tanh\Big(\tfrac{1}{|N(v)|}\textstyle\sum_{u\in N(v)} h_u \; W\Big), \qquad v\in A_t,
\end{equation}
with $F$-dimensional states ($F=64$). Recomputing $v$ costs $w_v=\deg(v)+F$ units of work
(aggregation plus dense transform), so the work of partition $P_i$ is
$W_i=\sum_{v\in P_i\cap A_t} w_v$, and it changes from window to window.

\paragraph{Execution model.}
The vertex set is split into $k$ disjoint partitions $P_1,\dots,P_k$
($\bigcup_i P_i=V_t$, $P_i\cap P_j=\emptyset$). Each partition is a bulk-synchronous worker
\cite{malewicz2010pregel} that owns its vertices and keeps read-only ghost copies of remote
neighbours. Partition $i$ has $\tau_i$ threads from a fixed budget $\sum_i\tau_i=T_{\mathrm{thr}}$.
A superstep consists of local computation, a barrier, and a halo exchange that refreshes the
ghosts of changed (active) vertices. Its time is governed by the slowest partition and by the
communication volume
\begin{equation}\label{eq:cv}
  \mathrm{CV} = \sum_{v} \big|\{P(u): u\in N(v)\}\setminus\{P(v)\}\big|,
\end{equation}
the number of ghost copies \cite{hendrickson2000models}. Unlike the number of ``boundary
edges'', which always equals twice the edge cut, CV counts each value once per receiving
partition.

\paragraph{Platforms.}
The executor runs on one multicore machine, and its time is always \emph{measured}. To study
deployments where partitions sit on different machines, we additionally emulate an
interconnect. The halo exchange of every superstep is charged
$\lambda\cdot\#\mathrm{msgs}+4F\cdot\mathrm{CV}_{A}/B$ (latency $\lambda$, bandwidth $B$,
$\mathrm{CV}_A$ = ghosts of active vertices), and vertex migration is charged at $B$. We use three
platforms: shared memory (\emph{SM}; migration at memory bandwidth, 10\,GB/s), emulated
10\,GbE ($B=1.25$\,GB/s, $\lambda=50\,\mu$s) and emulated 1\,GbE ($B=125$\,MB/s, $\lambda=100\,\mu$s).

\paragraph{Actions.}
Before executing window $t$, the runtime applies one action $a_t\in\mathcal{A}$:
\begin{itemize}
  \item \NOA{} keeps the partitioning and the thread allocation.
  \item \TA{} reallocates threads in proportion to the current work, $\tau_i\propto W_i$, with
        $\tau_i\ge 1$ and $\sum_i \tau_i=T_{\mathrm{thr}}$. It costs about 1\,ms and migrates nothing.
  \item \RB{} runs ACGR (Sect.~\ref{sec:acgr}), a local, communication-aware refinement that
        migrates few vertices, and then re-provisions the threads.
  \item \RP{} runs METIS \cite{karypis1998metis} from scratch, relabels the new parts to maximise
        their overlap with the old ones (scratch-remap \cite{schloegel2000unified}), and
        re-provisions the threads. It gives the best quality but migrates many vertices.
\end{itemize}
New vertices that arrive between decisions are placed online with LDG \cite{stanton2012ldg},
so an unmaintained partitioning degrades over time.

\paragraph{Objective.}
The goal is to minimise the end-to-end time of the stream,
\begin{equation}\label{eq:J}
  J = \sum_{t=1}^{T} \Big( T^{\mathrm{exec}}_t + T^{\mathrm{net}}_t
      + T^{\mathrm{adapt}}(a_t) + T^{\mathrm{mig}}(a_t) + T^{\mathrm{ctrl}}_t \Big),
\end{equation}
where $T^{\mathrm{exec}}$ is the measured superstep time, $T^{\mathrm{net}}$ the emulated network
time (zero on SM), $T^{\mathrm{adapt}}$ the measured run time of ACGR or METIS,
$T^{\mathrm{mig}}$ the migration time, and $T^{\mathrm{ctrl}}$ the controller's inference time.
Every term is a cost in seconds and is minimised. This replaces the weighted reward of our
earlier version, which mixed rewarded and penalised terms of different units with hand-set
weights. Partition quality (cut, CV, balance) affects $J$ only through the time it causes, so
no weights $w_1,\dots,w_4$ have to be chosen. The trade-off the controller must learn is
explicit in~\eqref{eq:J}: adapting costs time now and saves time later.
